\chapter{System Design}
\label{ch:system-design}

\section{System Overview}
The system is designed as a modular pipeline that learns normal host behavior and
flags abnormal behavior without relying on pre-labeled attacks. It continuously
collects host activity, converts it into behavioral graphs, trains a graph
autoencoder on normal graphs, and then scores new graphs for anomalies. The design
emphasizes clear boundaries between modules so that data collection, graph
construction, learning, detection, and evaluation can be improved independently.

\section{System Architecture}
Figure~\ref{fig:system-architecture} presents the end-to-end architecture. It
shows how raw events flow into graph construction, how the model learns normal
patterns, and how alerts are generated and visualized. The diagram also highlights
the evaluation loop where simulated attacks are injected to test detection.

The figure summarizes the full pipeline and the feedback loop that enables
evaluation of detection performance.

\begin{figure}[htbp]
	\centering
	\includegraphics[angle=90,origin=c,height=0.55\textheight,keepaspectratio]{\detokenize{architecture.png}}
	\caption{Overall system architecture for the zero-day detection pipeline.}
	\label{fig:system-architecture}
\end{figure}

This architecture view provides the reference flow for the module descriptions
that follow.

\section{Module Descriptions}
\subsection{System Monitoring and Data Collection}
The data collection module is responsible for observing real host activity. It
captures process actions, file access, and network activity at fixed intervals
and converts them into a common event format. The output is a clean stream of
normalized events that can be consumed by the graph builder.

The figure shows the flow from system sources through the collector into a
normalized event stream.

\begin{figure}[htbp]
	\centering
	\includegraphics[angle=90,origin=c,height=0.55\textheight,keepaspectratio]{\detokenize{system monitoring module.png}}
	\caption{System monitoring pipeline that produces normalized raw events.}
	\label{fig:system-monitoring}
\end{figure}

This module establishes the raw event stream used by all downstream stages.

\subsection{Behavioral Graph Construction}
The graph construction module groups events into time windows and converts each
window into a behavioral graph. Nodes represent processes, files, and sockets.
Edges represent interactions such as executes, reads, writes, spawns, and
connects. Node features are derived from type and connectivity statistics so the
model can learn structure as well as activity intensity.

The figure outlines how events are grouped into windows and converted into nodes,
edges, and feature vectors.

\begin{figure}[htbp]
	\centering
	\includegraphics[width=0.95\textwidth,height=0.68\textheight,keepaspectratio]{\detokenize{behavioural graph constructor.png}}
	\caption{Behavioral graph construction workflow with preprocessing and feature creation.}
	\label{fig:graph-construction}
\end{figure}

It shows the transformation from events to structured graphs used for learning.

\subsection{Graph Autoencoder Training and Baseline}
The learning module trains a graph autoencoder only on normal graphs. The encoder
compresses each graph into a latent representation using GCN layers, while the
decoder reconstructs the adjacency structure. Reconstruction error becomes the
anomaly signal. The training output includes the mean and standard deviation of
errors, which define the baseline threshold for detection.

The figure highlights the encoder--decoder flow and how the baseline statistics
are derived from reconstruction error.

\begin{figure}[htbp]
	\centering
	\includegraphics[width=0.95\textwidth,height=0.68\textheight,keepaspectratio]{\detokenize{autoencoder and learner.png}}
	\caption{Graph autoencoder training pipeline and baseline statistics computation.}
	\label{fig:autoencoder-pipeline}
\end{figure}

This diagram clarifies how the baseline threshold is computed from normal data.

\subsection{Detection and Alerting}
During detection, each new graph is passed through the trained autoencoder to
compute a reconstruction score. If the score exceeds the learned threshold, the
graph is flagged as anomalous and an alert is created with severity information.
This module also records detection results for reporting and analysis.

\subsection{Attack Simulation and Evaluation}
To validate detection behavior, the system includes an attack simulator that
generates synthetic streams for common attack categories. These events are merged
with normal activity and converted to graphs so the detector can be evaluated
under controlled conditions.

The figure shows how normal and simulated attack streams are combined before
graph conversion and scoring.

\begin{figure}[htbp]
	\centering
	\includegraphics[width=0.95\textwidth,height=0.68\textheight,keepaspectratio]{\detokenize{attack simulator.png}}
	\caption{Attack simulation pipeline for controlled evaluation.}
	\label{fig:attack-simulator}
\end{figure}

The simulated streams ensure repeatable and measurable detection tests.

\subsection{Visualization and Reporting}
The visualization module provides graph views and summary plots that allow a
reviewer to compare normal and abnormal behavior. These views support qualitative
validation of detection outcomes and help explain why a graph was flagged.

\section{End-to-End Data Flow}
The system executes the following sequence:
\begin{enumerate}
	\item Capture raw system events during normal operation.
	\item Convert events into time-windowed behavioral graphs.
	\item Train the graph autoencoder on normal graphs only.
	\item Compute baseline statistics for anomaly thresholds.
	\item Score new graphs and generate alerts if thresholds are exceeded.
	\item Inject simulated attacks and evaluate detection performance.
	\item Export logs and visual summaries for reporting.
\end{enumerate}

\section{Summary}
This chapter redefines the system design around the major module diagrams. Each
figure corresponds to a pipeline stage, from event capture through graph learning
to detection and evaluation, providing a consistent and visually grounded
description of the full system.
